%
% Settings file
%
\include{settings}	% Document settings file

%
% We use the book class
%
\documentclass[12pt]{book}

% toggle different modes
% \Drafttrue
% \DownscaledFinalDoctrue
\CompactLayouttrue %smaller chapter/section headers and less whitespace around them 

%
% The geometry package allows consistent control of page layout
%
\usepackage[paper = a4paper,         %To be shrunk to 80% when printed
			twoside,                 %Two-side mode, switches margins on 
			bindingoffset = 2mm,     %Offset for binding side of page
			hmargin = 25mm,          %Left and right margin
			vmargin = 25mm,          %Top and bottom margin
			dvipdfm]
			{geometry}


%
% Font settings
%
\usepackage[T1]{fontenc}			% Activate Type 1 fonts
%\usepackage{mathptmx}       			% Use Times font, also for math
\usepackage[scaled]{helvet} 			% for sans serif fonts (\textsf{...} or \sffamiliy)
%\usepackage[scaled]{luximono}			% for monospaced fonts (\texttt{...} or \ttfamily)

% 
% Header settings
%
\usepackage{sectsty}        			% Change section and chapter header
%\allsectionsfont{\usefont{T1}{phv}{bc}{n}\selectfont} % Set chapter/section header to narrow helvectica (arial-like)
\usepackage{titlesec} 				%fix header fonts
\ifCompactLayout
    % more compact chapter and section titles (optional, article style)
    \titleformat{\chapter}
      {\normalfont\sffamily\Large\bfseries}
      {\thechapter}{10pt}{}
      \titlespacing*{\chapter}{0pt}{0pt}{20pt}
    \titleformat{\section}
      {\normalfont\sffamily\large\bfseries}
      {\thesection}{10pt}{}
    \titleformat{\subsection}
      {\normalfont\sffamily\bfseries}
      {\thesubsection}{10pt}{}
    \titleformat{\subsubsection}
      {\normalfont\sffamily\bfseries}
      {\thesubsubsection}{10pt}{}
\else
    % less compact chapter and section titles (default, book style) 
    \titleformat{\title}[display]
      {\normalfont\rmfamily\huge\bfseries\color{black}}
      {\chaptertitlename\ \thechapter}{20pt}{\Huge}
    \titleformat{\chapter}[display]
      {\normalfont\rmfamily\huge\bfseries\color{black}}
      {\chaptertitlename\ \thechapter}{20pt}{\Huge}
    \titleformat{\section}
      {\normalfont\rmfamily\Large\bfseries\color{black}}
      {\thesection}{1em}{}
    \titleformat{\subsection}
      {\normalfont\rmfamily\Large\bfseries\color{black}}
      {\thesubsection}{1em}{}
\fi

%
% Set fancy page header using fancyhdr package
% 
\usepackage{fancyhdr}
\pagestyle{fancy}

% Ensure that the chapter and section headings are in lowercase
\setlength{\headheight}{15pt}
\renewcommand{\chaptermark}[1]{\markboth{#1}{}}
\renewcommand{\sectionmark}[1]{\markright{\thesection\ #1}}

% Delete current section for header/footer
\fancyhf{}

% Define header/footer layout
\fancyhead[LE,RO]{\bfseries\thepage}
\fancyhead[LO]{\bfseries\rightmark}
\fancyhead[RE]{\bfseries\leftmark}
\renewcommand{\headrulewidth}{0.5pt}

% make space for the rule
\fancypagestyle{plain}{
	\fancyhead{} 				%get rid of the headers on plain pages
	\renewcommand{\headrulewidth}{0pt} 	% and the line
}

%
% Set header height to 30pt to make room for two rows in the header
% (useful if you have very long chapter names)
%
%\headheight 30pt

% \raggedbottom					% may be used to avoid varying line spaceing between paragraphs at the cost of the last line on each page not appearing at the same height.

%
% Packages we may need (can be disabled/enabled depending on your requirments)
%
\usepackage{graphicx}       			% Handles figures
\usepackage[utf8]{inputenc} 			% We want æøå
\usepackage{amsmath}
\usepackage{amsfonts}
\usepackage[mathcal]{euscript} 			% For caligraphy fonts
\usepackage{setspace}			       	% Easy setting of line spacing
%\usepackage{forloop}				    % For-loops!
%\usepackage[numbers,sort&compress]{natbib}     % Sort numerical keys for multiple cites
\usepackage[authoryear,round]{natbib}	% for author year citations
\setcitestyle{nosort}				    % preserve order of citations in body text (this may collide with some bibliography styles) 
%\usepackage[nosort]{hypernat}       	% Avoid breaking 'backref' option of hyperref package
\usepackage{float}
\usepackage[dvipsnames]{color}   		% named colors
\usepackage[procnames]{listings} 		% beautiful listings
\usepackage[font={small,it,rm}]{caption}
\usepackage[final]{pdfpages}			% to include your papers as pdf
\usepackage{epigraph}				    % for quotation
\setlength\epigraphwidth{8cm}
%\usepackage{epstopdf}
\usepackage{array}				        % to align equations, figures etc.
\usepackage{booktabs}       			% Publication quality tables
\usepackage{tabularx}				    % for advanced table layout
\usepackage[table]{xcolor}
\usepackage{tikz}				        % include after xclolor!!!
%\usepackage{multirow}
% \usepackage{units}				    % semantically represent numbers with units
\usepackage{siunitx}				    %siunitx may be the better alternative for numbers with units
	\sisetup{separate-uncertainty}		% define how to display uncertainties
	\newcommand{\SIh}[2]{\SI[number-unit-product={\text{-}}]{#1}{#2}} % command to produce hyphen separated numbers and units as in "100-m level"
\usepackage{fancytabs}                  % required for creating thumb-indices for the papers (see example in 5Papers.tex)
\usepackage[nomessages]{fp}

\usepackage[euler]{textgreek}
\usepackage{upgreek}
\usepackage{pdflscape}
\usepackage{wrapfig}				% for figures embedded in text
\usepackage{textcomp}
\usepackage{sidecap} 				% for figure caption on the side
\usepackage{enumitem} 				% for possibility to adjust item left margin indentation
\usepackage[hyphens,spaces,obeyspaces]{url}
\usepackage[multiple]{footmisc}			% 
% \usepackage{atveryend}              % for execution of code at the end of compilation process
%\usepackage{thumbs}

% useful packages while working on the draft
\usepackage[textsize=scriptsize]{todonotes}	% add notes to your pdf
\usepackage{soul}
\usepackage[normalem]{ulem}			% provides strike-out command \sout{}
\usepackage{lipsum}                 % create dummy text

\usepackage[
%	dvipdfmx,        			% Use dvipdf driver, disable for overleaf
	backref = page,   		      	% List citing occurences in the References, option [section] (default) doesn't look nice - [page] works better, can also be disabled 
	colorlinks,      			% Colored links
	citecolor=black,  			% Color of cite links
	linkcolor=black,  			% Color of links
	urlcolor=blue,   			% Color of urls
	]{hyperref}
%\newcommand{\aref}[1]{\autoref*{#1}} 		% Prevent links
\newcommand{\aref}[1]{\autoref{#1}} 		% Enable links



%
% Include user-defined macros
%
%====================================================
%-------------Macro definitions go here--------------
%====================================================

%
% Differentials
%
\newcommand{\tdiff}[2]{\ensuremath{\frac{d#2}{d{#1}}}}
\newcommand{\tdifforder}[3]{\ensuremath{\frac{d^{#2}#3}{d{#1}^{#2}}}}
\newcommand{\pdiff}[2]{\ensuremath{\frac{\partial#2 }{\partial#1}}}
\newcommand{\pdifforder}[3]{\ensuremath{\frac{\partial^{#2}#3}{\partial{#1}^{#2}}}}

%
% Linear algebra
%
\renewcommand{\vec}[1]{\ensuremath{\mathbf{#1}}}
\newcommand{\mat}[1]{\ensuremath{\mathbf{#1}}}
\newcommand{\tildemat}[1]{\ensuremath{\widetilde{\mat{#1}}}}


%
% Misc macros
%
%\renewcommand{\equiv}[0]{\ensuremath{:=}}
\newcommand{\etal}{\textit{et al. }}
\newcommand{\paperheader}[2]{\noindent\textbf{Paper #1}: \textit{#2}\\}
% \newcommand{\paperitem}[3]{\noindent\textbf{Paper #1}: \textit{#2}\vspace{1em}\\\noindent #3\vspace{2em}}
\newcommand{\paperitem}[4]{\noindent\textbf{Paper #1}: \textbf{#2}\vspace{1em}\\\textit{#3}\vspace{1em}\\\noindent #4\vspace{2em}}
\newcommand{\tfinal}{\ensuremath{T_{\text{f}}}}
\newcommand{\papernum}[1]{\textbf{#1}}

\newcommand{\footurl}[1]{\footnote{\url{#1}}}		% url as footnote

% some shortcuts for often used units (package siunitx)
\DeclareSIUnit\met{\metre}
\DeclareSIUnit\km{\kilo\metre}
\DeclareSIUnit\sqkm{\square\kilo\metre}
\DeclareSIUnit\sec{\second}
\DeclareSIUnit\min{\minute}
\DeclareSIUnit\kg{\kilo\gramm}
\DeclareSIUnit\hPa{\hecto\pascal}
\DeclareSIUnit\mps{\met\per\sec}
\DeclareSIUnit\wpsm{\watt\per\square\met}

% default figure width
\newcommand{\figw}{0.9\linewidth}

% referencing sections, figures, equations, tables
\newcommand{\eref}[1]{~(\ref{#1})}
\newcommand{\Fref}[1]{Figure~\ref{#1}}
\newcommand{\Sref}[1]{Section~\ref{#1}}
%\newcommand{\sref}[1]{Sect.~\ref{#1}}
\newcommand{\Tref}[1]{Table~\ref{#1}}
\newcommand{\Eref}[1]{Equation~\ref{#1}}     

% shortcuts for often used variables
%\newcommand{\dTdz}{\frac{\partial \theta}{\partial z}}
%\newcommand{\ldTdz}{\partial \theta /\partial z}
%\newcommand{\wTv}{\overline{w'\theta_v '}}

% shortcuts for institutions
\newcommand{\GFI}{Geophysical Institute, University of Bergen}

% colored text
\newcommand{\tcr}[1]{\textcolor{red}{#1}}

% papers
\newcommand{\PI}{\mbox{Paper~I}}
\newcommand{\PIT}{Title of \PI}
\newcommand{\PIA}{Authors of \PI}
\newcommand{\PIJ}{\noindent \textit{Journal}, \textbf{Volume}, Pages (year), doi:10.xxx/xxx}


%
% Theorem environments
%
\newtheorem{definition}{Definition}


%
% Counters
%
\newcounter{ct}

%
% For loop to include papers
%
\newcommand{\includepaperpages}[2]
{
	\forloop{ct}{0}{\value{ct} < #2}
	{
		\begin{figure}[ht]
			\includegraphics[width=.99\textwidth]{#1\the\value{ct}.ps}
		\end{figure}
		\clearpage
	}
}


%
% TeX is very proud of its hyphenation engine, to the point where it 
% hyphenates _everything_ just to show off.  Increasing penalty for 
% hyphenation, and increase tolerance for underfull boxes can make 
% the text easier to read, at the expense of making the text less aligned
% to the right margin
%
\hyphenpenalty=10
\tolerance=1000



\title{{\fontsize{25}{24}\sffamily\textbf{Thesis Title}}
	\author{
	\large\textbf{Your Name}\vspace{4em}\\
		\includegraphics[width=74mm]{figures/uglo}\vspace{4em}\\
		Dissertation for the degree of Philosophiae Doctor (PhD)\\ at the University of Bergen}
        %\vspace{3.5em}\\
		%Geophysical Institute, \\
		%University of Bergen
	%}
	\huge\date{2020}
}


%====================================================
%------------------ BEGIN DOCUMENT ------------------
%====================================================
\begin{document}
%Cause all references in bibtex file to appear in the 'References' section, even
%if they are not explicitly cite'ed in the document
%\nocite{*}

%Set the fontsize and baselineskip, if something other than 10 or 12 is required
\ifDownscaledFinalDoc
	\fontsize{\TextSize}{\BaseLineSkip}
	\selectfont
\fi

%Double line spacing for draft
\ifDraft
	\doublespacing
\fi

\renewcommand{\familydefault}{\sfdefault} %KD

%--------------------------------------------------------------------------------------------------
% FRONT MATTER
%--------------------------------------------------------------------------------------------------
\pagenumbering{roman}
\clearpage\maketitle
\thispagestyle{empty}
\frontmatter
\normalfont\rmfamily
\include{01scientific_environment}
\chapter{Acknowledgements}
\ldots
% If you want to acknowledge the people creating this template you can thank Kjersti Birkeland Daae and Stephan Kral for making it available on overleaf and Tore Birkeland and Raymond Nepstad for creating the original version of the template. 
%\vspace{5em}
\vspace*{\fill}
\begin {flushright}
Your Name\\
Place, Date
\end{flushright}
\clearpage

%quote
\vspace*{\fill}
\epigraph{"quotation"}{\textit{Reference}}

\include{03abstract}
\chapter{Outline}

This thesis consists of an introductory part and four scientific papers. 
Chapter \ref{chap:intro} gives an introduction to \ldots
Objectives, data and methods, as well as a detailed description of \ldots in chapter \ref{chap:thisstudy}. 
A brief summary of the papers is given in chapter \ref{chap:intropaper}\ldots
The papers included in this thesis (chapter \ref{chap:papers}) are:

\begin{enumerate}
\item authors, (year) \textit{title}, journal \textbf{Volume/Number} 

\item authors, (year) \textit{title}, journal \textbf{Volume/Number} (Accepted, Aug 2019)

\end{enumerate}

\tableofcontents
%\listoffigures
%\lstlistoflistings

%--------------------------------------------------------------------------------------------------
% MAIN MATTER
%--------------------------------------------------------------------------------------------------
\mainmatter

%
% Main chapters
%
\chapter{Introduction}
\label{chap:intro}
\ldots

\begin{figure}[htbp]
\centering
	\includegraphics[width=12cm]{figures/uglo} 
	\caption{caption describing the figure}
	\label{fig:1}
\end{figure}
%%----------------------------------------------------------------------
%%----------------------------------------------------------------------
\section{sectiontitle}
\label{sec:i1}
\lipsum[1]
\begin{wrapfigure}{r}{0.5\textwidth}
    \centering
    \includegraphics[width=.47\textwidth]{figures/uglo.eps}
    \caption{wrapped figures may need careful adjustments to ensure that they don't affect the layout of the following page.}
    \label{fig:wf}
\end{wrapfigure}
\lipsum[1]

\ldots
%%----------------------------------------------------------------------
%%----------------------------------------------------------------------
         % Introduction
%% here you should identify the knowledge gaps you fill with this thesis, formulate your research questions and hypothesis etc. 
\chapter{This study}
\label{chap:thisstudy}
\section{Objectives}
Processes related to \ldots
\subsection{\ldots}
\ldots 

\begin{table}[!htbp]\centering
\def\arraystretch{1.3}
\caption{tablecaption}
\vspace{5pt}
    \arrayrulecolor{black}
    \small
\setlength{\tabcolsep}{4pt}
\begin{tabular}{lcrccccccll}
\hline
 c1 & c2 & c3 & c4 & c5 & c6 & c7 \\
\end{tabular}
\label{Tab:1}
\end{table}
\normalsize

\newpage
\begin{landscape}
\begin{table}[!htbp]\centering
\def\arraystretch{1.3}
\caption{tablecaption}
\vspace{5pt}
    \arrayrulecolor{black}
    \small
\setlength{\tabcolsep}{4pt}
\begin{tabular}{lcrccccccll}
\hline
 c1 & c2 & c3 & c4 & c5 & c6 & c7 \\
\end{tabular}
\label{Tab:2}
\end{table}
\end{landscape}
\normalsize


\chapter{Introduction to the papers}
\label{chap:intropaper}
%
% Paper I
%
\paperitem{I}
{\PIT}
{authors (year), journal, volume(number)}
{{\PI} investigates \ldots}

  % Introduction to the papers
\include{4Summary_Conclusions}
\chapter{Scientific results}
\label{chap:papers}


%% some additional settings and commands relevant for this section
% command to include watermark, e.g. for accepted papers
\newcommand{\watermark}[1]{\tikz[remember picture, overlay] \node [red,opacity=0.3,rotate=45,scale=5.0] at (current page.center) {#1};}

% commands and settings for thumb indices created with \fancytab command
\newcommand{\nop}{3}                                                % put in the number of papers you have in your thesis
\newcommand{\rtw}{.12}                                              % relative width of thumb index (0.18 is used by printing company)
\newcommand{\thumbheight}{\dimexpr\paperheight/\nop\relax}
\FPeval{\rmpw}{1-\rtw}
\fancytabsStyle{\Huge\scshape}
\fancytabsHeight{\thumbheight}
\fancytabsWidth{\rtw\paperwidth}
\fancytabsCount{\nop}
\fancytabsLeftColor{black}
\fancytabsRightColor{black} 
\fancytabsTop{0mm}
% \fancytabsTextVPos{0.5} 
% \fancytabsTextHPos
% \fancytabsGap 
%\fancytabsFloor 
\fancytabsRotate{90}                    % change to 0 for non rotated index 
\newcommand{\thumbindex}[2]{% 
    \let\eastwest\undefined 
    \let\multiplier\undefined 
    \fancytab[RIGHT]{\scalebox{1.5}{\textbf{\textcolor{white}{#1}}}}{#2}
    }
% \thumbindex requires two inputs, text and a position index (numer between 1 and \nop (number of papers)

%!!! note that the printing company can also add the thumbindexes, but you have to ask for it specifically, and you should create enough space (change the right margin as done here using the minipage environment). You may also want to move the text slightly to the left, since the printing company "cuts" 3 mm away in the resized 17 cm x 24 cm format (for A4 format this corresponds to about 3.7 mm)

% note that the fancytabs-package is actually meant to be used in the definition of a new pagestyle (\fancypagestyle). One could define such a page style for this section but I don't feel its worth the effort.

%--------------------------------------------------------------------------------------------------

\pagestyle{empty}
% Paper I
%--------------------------------------------------------------------------------------------------
\cleardoublepage
\chapter*{Paper I}
\begin{minipage}[c]{\rmpw\linewidth}
\thumbindex{I}{1}
\section*{title}
\noindent author1 author2 author3 author4 author5 author6 author7 author8 author9 author10 author11 author12\\
\noindent \textit{journal}, \textbf{volume/number} (year)
\end{minipage}
\cleardoublepage
%\includepdf[pages=-]{papers/paper-1.pdf} %KD
%\includepdf[pages=1,pagecommand={\watermark{ACCEPTED}}]{papers/paper-1.pdf}
%\includepdf[pages=2-]{papers/paper-1.pdf}

% Paper II
%--------------------------------------------------------------------------------------------------
\chapter*{Paper II}
\begin{minipage}[c]{\rmpw\linewidth}
\thumbindex{II}{2}
\section*{title}
\noindent authors\\
\noindent \textit{journal}, \textbf{volume/number} (year)
\end{minipage}
\cleardoublepage
%\includepdf[pages=-]{papers/paper-2.pdf} %KD

% Paper III
%--------------------------------------------------------------------------------------------------
\chapter*{Paper III}
\begin{minipage}[c]{\rmpw \linewidth}
\thumbindex{III}{3}
\section*{title}
\noindent authors\\
\noindent \textit{journal}, \textbf{volume/number} (year)
\end{minipage}

% consider adding a copyright statement for unpublished work or non-open source publications (you have to ask for permission by the journal if you want to publish your thesis, e.g. in BORA) 
\newpage
\vspace*{\fill}
\noindent\textcopyright~The copyright holder for this preprint is the authors. All rights reserved. No further reuse allowed without permission. 

%\includepdf[pages=-]{papers/paper-3.pdf} %KD

%--------------------------------------------------------------------------------------------------
               % The papers, with cover pages
%
% Appendices
%
%\appendix
%\include{appendixA}             % Appendix A: Something


%--------------------------------------------------------------------------------------------------
% BACK MATTER
%--------------------------------------------------------------------------------------------------
\backmatter

%
% The bibliography
%

%
% Physics-style bibliography
%
%\bibliographystyle{iopart-num}

%
% Modified mathematics bibliography. Uses acm.bst
%

\bibliographystyle{agufull08}

\bibliography{Thesis}

\end{document}

